% Additive reader entry; no original source chapter is replaced.
\chapter*{Reader's laboratory: from equations to executable evidence}
\addcontentsline{toc}{chapter}{Reader's laboratory: from equations to executable evidence}
\noindent Mathematics for visual computing begins before we draw the first
coordinate axis. A physical observation, a mathematical model, a finite
representation and a software implementation are different objects.
Confusing them can make a program reproduce its own mistakes while apparently
passing its tests. This book will move through those objects explicitly.

Suppose a sensor reports a point on a surface. The physical point has a
location and units. The mathematical point belongs to an affine space;
subtracting two positions gives a displacement, whereas adding two positions
without a defined affine convention has no geometric meaning. A coordinate
frame expresses that point as numbers, but those numbers depend on the basis
and origin. Floating-point storage adds rounding; the algorithm adds
approximation, stopping thresholds and possibly statistical uncertainty.
A rendered image of the result adds another interpretation layer. Each
transition must have a declared contract.

\paragraph{One small example.}
Take a local frame with origin \(o=(5,-1)\) and world-space axes
\(a=(2,0)\) and \(b=(1,3)\). A point with local coefficients
\(q=(2,-1)\) maps to
\[
p=o+2a-b=(5,-1)+(4,0)-(1,3)=(8,-4).
\]
A result of \((8,-4)\) is a checkable computation; it is not yet proof
that the inverse frame transformation handles arbitrary inputs. The
basis determinant is \(2\cdot3-0\cdot1=6\), so this particular
basis is invertible. Replacing \(b\) by \((1,10^{-15})\) can create a
poorly conditioned inverse even though the exact determinant is nonzero.
We should therefore check mathematical existence, numerical sensitivity
and implementation failure paths separately.

\paragraph{The five questions to ask at every new algorithm.}
First: what object and quantity is being represented, in what coordinate
system and units? Second: which assumptions make the formula valid?
Third: what is the error source---measurement, discretization, rounding,
conditioning or an incorrect model? Fourth: what independently checkable
result or invariant should be recovered? Fifth: what counterexample would
falsify the implementation's apparent success?

\paragraph{Our testing convention.}
Write down input geometry, reference answer, numerical tolerance and
failure policy before running an experiment. Use identical random seeds
only for repeatability, not as evidence of statistical independence.
When possible, test a second method: check a numerical integral with an
antiderivative, an optimization gradient with finite differences, a linear
solve by residual and conditioning, or a heat solver against a known
solution. An exact analytical image, a simulation result and a measured
benchmark must be labeled differently.

\paragraph{Cumulative workbench.}
Our executable companion begins with a 2D coordinate-frame witness,
adds stable solves and derivatives, grows into an uncertainty-and-fitting
laboratory, and ends in a field solver and geometric atlas. The reader
should keep both the code and a plain-language laboratory log recording
what was derived, what was tested and what remains uncertain.
